% =====================================================================
% Anhang C: Datenbank-Schema (DDL)
% =====================================================================

\chapter{Datenbank-Schema (DDL)}

\label{app:schema}

Dieser Anhang enthält das vollständige SQL-Schema, wie es von
Alembic-Migration~0001 erzeugt wird.

\section{Stammdaten}

\begin{lstlisting}[language=sql, caption={Stammdaten-Tabellen},label=lst:ddl-stamm]
CREATE TABLE venues (
    venue_id   SERIAL PRIMARY KEY,
    code       TEXT UNIQUE NOT NULL,
    name       TEXT,
    timezone   TEXT,
    country    CHAR(2)
);

CREATE TABLE vendors (
    vendor_id   SERIAL PRIMARY KEY,
    code        TEXT UNIQUE,
    license     TEXT,
    homepage    TEXT
);

CREATE TABLE instruments (
    instrument_id   BIGSERIAL PRIMARY KEY,
    asset_class     TEXT NOT NULL,
    name            TEXT NOT NULL,
    currency        CHAR(3),
    primary_venue   INTEGER REFERENCES venues(venue_id),
    is_active       BOOLEAN DEFAULT TRUE,
    created_at      TIMESTAMPTZ DEFAULT now()
);

CREATE TABLE instrument_identifiers (
    instrument_id  BIGINT REFERENCES instruments(instrument_id),
    scheme         TEXT NOT NULL,
    value          TEXT NOT NULL,
    PRIMARY KEY (scheme, value)
);
CREATE INDEX ON instrument_identifiers(instrument_id);

CREATE TABLE trading_calendars (
    venue_id        INTEGER REFERENCES venues(venue_id),
    date            DATE,
    is_trading_day  BOOLEAN,
    PRIMARY KEY (venue_id, date)
);
\end{lstlisting}

\section{Preis-Tabellen}

\begin{lstlisting}[language=sql, caption={Preis-Tabellen},label=lst:ddl-preise]
CREATE TABLE prices_raw (
    vendor_id      INTEGER REFERENCES vendors(vendor_id),
    vendor_symbol  TEXT NOT NULL,
    date           DATE NOT NULL,
    payload        JSONB NOT NULL,
    ingested_at    TIMESTAMPTZ DEFAULT now(),
    PRIMARY KEY (vendor_id, vendor_symbol, date)
);

CREATE TABLE prices_daily (
    instrument_id   BIGINT REFERENCES instruments(instrument_id),
    date            DATE NOT NULL,
    open            NUMERIC(18,8),
    high            NUMERIC(18,8),
    low             NUMERIC(18,8),
    close           NUMERIC(18,8),
    adjusted_close  NUMERIC(18,8),
    volume          BIGINT,
    vendor_id       INTEGER REFERENCES vendors(vendor_id),
    PRIMARY KEY (instrument_id, date)
);
CREATE INDEX ix_prices_daily_date ON prices_daily(date);
CREATE INDEX ix_prices_daily_instrument_id_date
    ON prices_daily(instrument_id, date);
\end{lstlisting}

\section{Makro-, FX- und Bond-Tabellen}

\begin{lstlisting}[language=sql, caption={Makro-, FX- und Bond-Tabellen},label=lst:ddl-makro]
CREATE TABLE corporate_actions (
    instrument_id  BIGINT REFERENCES instruments(instrument_id),
    ex_date        DATE,
    action_type    TEXT,
    value          NUMERIC(18,8),
    PRIMARY KEY (instrument_id, ex_date, action_type)
);

CREATE TABLE fx_rates_daily (
    ccy_from  CHAR(3),
    ccy_to    CHAR(3),
    date      DATE,
    rate      NUMERIC(18,8),
    vendor_id INTEGER REFERENCES vendors(vendor_id),
    PRIMARY KEY (ccy_from, ccy_to, date)
);

CREATE TABLE bond_yields (
    instrument_id  BIGINT REFERENCES instruments(instrument_id),
    date           DATE,
    yield          NUMERIC(10,6),
    price          NUMERIC(10,6),
    duration       NUMERIC(10,4),
    rating         TEXT,
    PRIMARY KEY (instrument_id, date)
);

CREATE TABLE macro_series (
    series_id  TEXT PRIMARY KEY,
    name       TEXT,
    source     TEXT,
    frequency  TEXT,
    unit       TEXT
);

CREATE TABLE macro_observations (
    series_id  TEXT REFERENCES macro_series(series_id),
    date       DATE,
    value      NUMERIC(18,6),
    PRIMARY KEY (series_id, date)
);
\end{lstlisting}

\section{DQ-Tabellen}

\begin{lstlisting}[language=sql, caption={DQ-Tabellen},label=lst:ddl-dq]
CREATE TABLE dq_runs (
    run_id    BIGSERIAL PRIMARY KEY,
    run_at    TIMESTAMPTZ DEFAULT now(),
    pipeline  TEXT,
    status    TEXT
);

CREATE TABLE dq_findings (
    run_id         BIGINT REFERENCES dq_runs(run_id),
    check_name     TEXT,
    severity       TEXT,
    affected_rows  INTEGER,
    details        JSONB,
    PRIMARY KEY (run_id, check_name)
);
\end{lstlisting}

\section{Indizes im Überblick}

\begin{table}[ht]
\centering
\small
\begin{tabular}{lll}
  \toprule
  \textbf{Tabelle} & \textbf{Index} & \textbf{Spalte(n)} \\
  \midrule
  \code{instruments}            & PK            & \code{instrument\_id} \\
  \code{venues}                  & UQ + PK       & \code{code} / \code{venue\_id} \\
  \code{vendors}                & UQ + PK       & \code{code} / \code{vendor\_id} \\
  \code{instrument\_identifiers}& PK            & (\code{scheme}, \code{value}) \\
                                  & \code{ix\_instrument\_identifiers} & \code{instrument\_id} \\
  \code{prices\_daily}          & PK            & (\code{instrument\_id}, \code{date}) \\
                                  & \code{ix\_prices\_daily\_date}        & \code{date} \\
                                  & \code{ix\_prices\_daily\_inst\_date}   & (\code{instrument\_id}, \code{date}) \\
  \code{macro\_observations}    & PK            & (\code{series\_id}, \code{date}) \\
  \code{dq\_findings}            & PK            & (\code{run\_id}, \code{check\_name}) \\
  \bottomrule
\end{tabular}
\caption{Alle Indizes des Schemas. \code{PK} = Primärschlüssel,
  \code{UQ} = Unique-Constraint.}
\end{table}
